\documentclass[11pt]{article}

\usepackage[margin=0.85in, top=0.75in, bottom=0.75in]{geometry}
\usepackage{hyperref}
\usepackage{parskip}

\setlength{\parindent}{0pt}
\setlength{\parskip}{8pt}

\pagestyle{empty}

\begin{document}

%====================
% Header
%====================
\begin{center}
    {\LARGE \textbf{Gautam Velpula}} \\[6pt]
    \href{mailto:gautamvelpula@gmail.com}{gautamvelpula@gmail.com} \;|\;
    \href{tel:6785390154}{678-539-0154} \;|\;
    \href{https://www.linkedin.com/in/gautamvelpula}{linkedin.com/in/gautamvelpula}
\end{center}

\vspace{14pt}

\today

\vspace{8pt}

Dear Hiring Team,

I'm writing to express my interest in the Head of Enterprise Architecture role at SiteOne Landscape Supply. This role describes exactly what I've been doing for the past four years---and what I'd like to do next at a company where enterprise architecture is positioned to have direct, measurable impact on business outcomes.

At GoTo Foods, I built the enterprise architecture function from the ground up. There was no EA charter, no governance model, no architecture standards when I arrived. I defined all of it: the EA service portfolio, decision rights, operating model, and the architecture principles and guardrails that now govern technology decisions across the enterprise. I secured executive buy-in by making EA practical and outcome-oriented---tied to revenue, cost, and risk---not theoretical frameworks that sit on a shelf.

The work that followed is directly relevant to what SiteOne needs:

\textbf{Technology modernization and portfolio rationalization.} I led the assessment of 7 independent brand technology platforms, defined the target-state architecture, and executed a consolidation that eliminated redundant systems, reduced technical debt, and cut time-to-market by 40\%. I defined the modernization implementation plan---migrating from monolithic legacy systems to microservices and API-led integration---and orchestrated delivery across business and IT teams.

\textbf{Connecting IT investment to business economics.} I don't present architecture decisions as technology choices. I frame them as business decisions---cost, revenue impact, risk, and time-to-value. I've led cross-functional capability mapping to identify improvement opportunities and recommended solution options that executives could evaluate on business terms, not just technical merit.

\textbf{Governance that enables, not restricts.} The EA governance model I built balances agility with compliance. Architecture standards and patterns are designed for reuse and adoption, not bureaucracy. Communities of practice and cross-functional partnerships ensure that guidance is integrated into execution, not imposed from above.

I'm also actively leading our enterprise's adoption of generative AI---defining the architectural patterns, guardrails, and governance that make it safe and scalable. This is an area where EA leadership matters enormously, and I'd bring both the strategic perspective and practical experience to SiteOne.

I'm based in Atlanta, which I understand is where this role is located. I'd welcome the opportunity to discuss how my experience building and operating an EA function aligns with SiteOne's vision.

\vspace{16pt}

Sincerely, \\
Gautam Velpula

\end{document}